\section{A reusable theorem for positive block transfer}
\label{sec:blocks}

Small-table optimization often changes the representation of control.
Two old states may become one new state if nearby tape bits retain the
distinction.  A head displacement or reversible tape change may also
be part of the representation.  Correctness must therefore compare
whole configurations, not just target-state numbers.

\begin{theorem}[Positive block transfer]
Let $M,N$ be deterministic machines with designated halt outcomes.
Let $I$ be an invariant of the old run and $F$ map its checkpoint
configurations into new configurations.  Suppose the initial
configurations correspond and halt/nonhalt is preserved by $F$.
For every nonhalting old checkpoint $C$ satisfying $I$, suppose there
are positive integers $a(C),b(C)$ such that an old block of $a(C)$
steps ends at a checkpoint $C'$, a new block of $b(C)$ steps ends
at $F(C')$, and neither block contains an interior halt.  Suppose
the blocks remain in their respective executions and the endpoint
again satisfies $I$.  Then the two initial runs halt together.
\end{theorem}

\begin{proof}
Concatenate the corresponding blocks.  Induction on their number
establishes corresponding endpoints.  If either run halts at an
endpoint, its corresponding endpoint is a halt; an interior halt is
excluded by hypothesis.  If all blocks continue, the two cumulative
clocks both tend to infinity because every increment is positive.
Consequently the concatenations cover every finite execution prefix
of both deterministic runs.  Both runs are then infinite.
\end{proof}

The theorem makes no assertion that unrelated row edits compose.
An edit can change the contexts entering another edited helper, so
the invariant, endpoint map, and interior paths must be checked for
the combined table.  A search depth limits the attempt to discover
blocks; after a finite family of universally guarded identities is
proved, it does not bound the simulated execution.

\subsection{The inherited 238-to-237 fold, proved on all tapes}

Let $M$ be the common 238-state input and $N$ the 237-state table,
both in inherited labels.  Delete 235, map it to 98, and make exactly
these retained-row changes~\cite{parent237}:
\begin{center}
\begin{tabular}{@{}lll@{}}
\toprule
Row & Old & New\\
\midrule
$98/0$ & $0R113$ & $0L240$\\
$139/1$ & $1R235$ & $1R98$\\
$156/1$ & $1R235$ & $1R98$\\
$240/1$ & $1R99$ & $1L87$\\
\bottomrule
\end{tabular}
\end{center}
The removed rows are $235/0=0R231$ and $235/1=1R16$.
All 470 other retained rows are unchanged.  Let
$F(q,h,T)=(\phi(q),h,T)$, where $\phi(235)=98$ and $\phi$
fixes every other state, including start 95 and halt 254.

The old-run invariant is
\begin{align*}
 q\in\{98,240\}&\Longrightarrow T(h-1)=0,\\
 q\in\{139,156\}&\Longrightarrow T(h-1)=1,\\
 q=235&\Longrightarrow T(h-2)=T(h-1)=1.
\end{align*}
Every incoming edge to 98 or 240 writes 0 and moves right.  The
only predecessors of 139 and 156 write 1 and move right.  Entry
to 235 is from 139 or 156, writes another 1, and moves right;
the predecessor invariant supplies its second left-hand 1.
The exact incoming rows are $86/0,187/0$ for 98;
$117/1$ for 139; $195/1$ for 156; $139/1,156/1$ for 235;
and $135/0,242/0$ for 240.  This proves induction closure.
At state 95 the invariant is vacuous for every tape and head position.

Only three old rows need longer new blocks:
\begin{center}\small
\begin{tabular}{@{}llll@{}}
\toprule
Old row & Guard & New state path & Steps\\
\midrule
$98/0$ & left 0 & $98\to240\to176\to113$ & 3\\
$235/0$ & left 11 & $98\to240\to87\to164\to118\to231$ & 5\\
$240/1$ & left 0 & $240\to87\to190\to99$ & 3\\
\bottomrule
\end{tabular}
\end{center}
Setting the initial head to zero for notation, the triples
$(\text{state},\text{position},\text{read bit})$ are
\begin{align*}
 &(98,0,0),(240,-1,0),(176,0,0);\\
 &(98,0,0),(240,-1,1),(87,-2,1),(164,-1,1),(118,0,0);\\
 &(240,0,1),(87,-1,0),(190,0,1).
\end{align*}
Every transition writes the bit it reads.  No cell outside $[-2,0]$
is touched, the complete tape is unchanged, and each block ends
one cell right at the old target.  All 473 other old rows commute
in one step through $F$, including $235/1$, which agrees with $98/1$.
The positive block theorem proves halting equivalence from state 95
on every binary tape.  If $a_t,b_t,c_t$ count old visits to the three
exceptional rows, the exact synchronized clock is
\[
             \tau(t)=t+2a_t+4b_t+2c_t,
             \qquad t\leq\tau(t)\leq5t.
\]

\subsection{A new independent finite checker}

The new \file{verification/verify_representation.py} begins with
every radius-$R$ binary window at old start state 95.  It computes
the least fixed point of the finite abstract transition relation,
branching both ways on each newly exposed exterior bit.  Induction
shows that this overapproximates every concrete old run on every
initial tape; it is not an empirical execution profile.

For every admitted checkpoint window, it simulates positive old
and new blocks and compares mapped whole-window endpoints.  It
rejects reads outside the represented window and any unavailable
representation-dependent cell.  Common arbitrary exterior tape
is untouched, so equality of finite overrides proves equality of
the complete tapes.  Halt stops a block, preventing interior halts.
Its optional state split excludes an old helper from checkpoints
only after verifying that all incoming edges to the helper come
from the deleted state or the helper itself.

On the inherited fold, radius 2, old cap 1 and new cap 5 give
3,282 cases: 3,266 blocks of lengths $(1,1)$, 12 of $(1,3)$, and
4 of $(1,5)$.  The same checker independently validates the earlier
context split.  A negative control restoring old helper row
$240/1=1R99$ is rejected at old state 235 with local window 11000.
That rejection concerns the proposed local certificate; it does
not claim the mutated blank run halts.
